\documentclass[11pt]{article}
% Shared preamble for all daily exercises
% Update this file to change settings (padding, parindent, etc.) for all exercises

\usepackage[margin=0.75in]{geometry}
\usepackage{amsmath}
\usepackage{amssymb}

% Custom settings
\setlength{\parindent}{0pt}
\setlength{\parskip}{0.2cm}

\title{Daily Exercise}
\author{Dhruman Gupta}
\date{29th March}

\begin{document}

% Common header for daily exercises
% This file can be included in other LaTeX documents

\maketitle

\noindent
\textbf{Course:} CS-3330, Theory of Computation

\vspace{0.2cm}

\noindent
\textbf{Question:}\noindent 
Consider the Post Correspondence Problem (PCP) over the binary alphabet $\{0, 1\}$. For any set $T$ of $n$ tiles, let $f(T)$ be the length of the shortest solution to the PCP instance $T$ if one exists, and $0$ otherwise. Define $g(n)$ to be the maximum value of $f(T)$ over all possible sets $T$ of exactly $n$ binary tiles. Intuitively, $g(n)$ measures how long the shortest solution can get in the worst case, when you are allowed to pick the worst possible $n$-tile binary PCP instance. Show that $g$ is not a computable function.

\textbf{Answer:}

We know that PCP over a binary alphabet is undecidable.

Suppose for the sake of contradiction that $g$ is computable. We will show that this lets us decide binary PCP.

Let $T$ be any PCP instance over $\{0,1\}$ with exactly $n$ tiles. Compute
\[
l = g(n).
\]
By definition of $g(n)$, if $T$ has any solution at all, then its shortest solution has length at most $l$.

Now enumerate all finite sequences of tiles of lengths $1,2,\dots,l$. For each such sequence, check whether it is a solution to $T$. This is a finite search, since there are only finitely many sequences of length at most $l$.

There are now two cases:
\begin{itemize}
    \item If during this search we find a solution, accept.
    \item If no solution is found after checking all sequences of length at most $l$, reject.
\end{itemize}

This procedure decides whether $T$ has a solution. Indeed, if $T$ has a solution, then by definition of $g(n)$ it has one of length at most $l$, so our search will find it. If no such sequence is found, then $T$ has no solution.

Hence binary PCP would be decidable. This is a contradiction, since binary PCP is undecidable.

Therefore $g$ is not computable.

\end{document}
